% Generated draft: requires mathematical and source review.
\documentclass[11pt]{article}
\usepackage[margin=25mm]{geometry}
\usepackage{fontspec}
\setmainfont{Latin Modern Roman}
\usepackage{amsmath,amssymb,mathtools,unicode-math}
\setmathfont{Latin Modern Math}
\usepackage{enumitem,needspace,xurl,hyperref}
\setlength{\parindent}{0pt}
\setlength{\parskip}{5pt}
\setlength{\emergencystretch}{3em}
\newcommand{\N}{\mathbb{N}}
\newcommand{\Z}{\mathbb{Z}}
\newcommand{\Q}{\mathbb{Q}}
\newcommand{\R}{\mathbb{R}}
\newcommand{\C}{\mathbb{C}}
\newcommand{\F}{\mathbb{F}}
\newcommand{\A}{\mathbb{A}}
\newcommand{\PP}{\mathbb{P}}
\newcommand{\E}{\mathbb{E}}
\newcommand{\eps}{\varepsilon}
\newcommand{\dd}{\,\mathrm{d}}
\begin{document}
\section*{UnsolvedMath ID: 700879}
\textbf{Source problem code:} catalogue entry
\subsection{Definitions and mathematical statement}
For a finite simple undirected graph \(\Gamma=(V,E)\), write
\[
N_\Gamma(v)=\{u\in V:\{u,v\}\in E\}
\]
for the open neighborhood of \(v\), and let \(\Gamma[N_\Gamma(v)]\) denote its induced graph. Under the usual nontrivial convention, \(\Gamma\) is strongly regular with parameters \((n,k,\lambda,\mu)\) if \(|V|=n\), \(0<k<n-1\), every vertex has degree \(k\), every two adjacent vertices have exactly \(\lambda\) common neighbors, and every two distinct nonadjacent vertices have exactly \(\mu\) common neighbors.

A finite generalized quadrangle of order \((s,t)\), with positive integers \(s,t\), is an incidence structure of points and lines in which each line contains \(s+1\) points, each point lies on \(t+1\) lines, two distinct points lie on at most one common line, and two distinct lines meet in at most one point. Two distinct points are collinear if they lie on a common line. For every point \(p\) and line \(L\) not incident with \(p\), there is exactly one point of \(L\) collinear with \(p\). Its point graph has the points as vertices, with distinct points adjacent precisely when they are collinear. The quadrangle is called thick when \(s,t\geq 2\).

Under the interpretation that a neighborhood is a generalized quadrangle when its induced graph is the point graph of such an incidence structure, the requested classification is:
\begin{quote}
Describe, up to graph isomorphism, all finite strongly regular graphs \(\Gamma\) such that, for every \(v\in V\), there exists a finite generalized quadrangle \(\mathcal Q_v\) satisfying
\[
\Gamma[N_\Gamma(v)]\cong \operatorname{Pt}(\mathcal Q_v),
\]
where \(\operatorname{Pt}(\mathcal Q_v)\) denotes its point graph.
\end{quote}
The source refers to Problem 12.58 for the intended meaning of generalized quadrangle, but that record is not supplied. Consequently, the point-graph interpretation, the finite nontrivial convention for strongly regular graphs, and whether thickness is required remain qualifications of this formulation. The displayed condition does not impose a common isomorphism type on the \(\mathcal Q_v\); the source does not explicitly require one. This is a statement-only draft, not a classification or a mathematical verification of the supplied open status.
\subsection{Short English statement}
Classify strongly regular graphs whose induced open neighborhoods are generalized quadrangle point graphs. This standard interpretation is conditional: the cited Problem 12.58 is missing, so the intended conventions, including thickness and possible degenerate cases, require review.
\subsection{Sources}
\begin{itemize}
\item Local UnsolvedMath registry, supplied complete source record: permanent ID \(700879\), exact source problem code \texttt{catalogue entry}, title \emph{Kourovka Notebook Problem 12.60}. The registry labels the problem open; this is supplied metadata, not an independently verified assessment. The supplied record was created and updated at \texttt{2026-10-16T14:00:00Z}, with upstream revision \texttt{d371ee917ac9647fbb06cd6fc529deec6f3dda82}.
\item \emph{The Kourovka Notebook (Unsolved Problems in Group Theory)}, Problem 12.60, attributed to A. A. Makhn\"ev, proposed in 1992 and associated with the 12th issue (1992). The supplied source also identifies the 21st issue edition (2026), \texttt{arXiv:1401.0300v48}. The registry describes the collection as Issues 1--21 (1965--2026), edited by E. I. Khukhro and V. D. Mazurov. These bibliographic details are supplied attribution and have not been independently checked.
\item The statement explicitly refers to Problem 12.58. Its contents are absent from the supplied record, so its definitions and qualifications could not be incorporated.
\end{itemize}
\end{document}
